\documentclass[10pt,a4paper]{amsart}
\usepackage[margin=19mm]{geometry}
\usepackage{amsmath,amssymb,mathtools}
\usepackage[scr=rsfs,cal=euler]{mathalfa}
\usepackage{xfrac}
\usepackage[unicode,hidelinks,bookmarks=false]{hyperref}
\newcommand{\ie}{i.e.\ }
\newcommand{\cf}{cf.\ }
\newcommand{\resp}{respectively\ }
\newcommand{\Subsection}{\subsection}
\providecommand{\nobreakdash}{\nobreak\hbox{-}}
\setlength{\emergencystretch}{2em}
\multlinegap0pt
\usepackage{enumerate}
\usepackage{amssymb}
\usepackage{float}
\usepackage{xcolor}
\usepackage{tikz}
\usepackage{algorithm}\usepackage{algpseudocode}
\newtheorem{lem}{Lemma}
\newtheorem{prop}{Proposition}
\title{Unimodular Equivalence of Integral Simplices}
\author{Feihu Liu$^{\color{blue} \dag}$, Sihao Tao$^{\color{blue} \S}$, and Guoce Xin$^{\color{blue} \P}$}
\date{Source: arXiv:2601.06819v1 (CC BY 4.0)}
\begin{document}
\maketitle

\noindent\emph{Source and licence.} Unimodular Equivalence of Integral Simplices. Version arXiv:2601.06819v1 (CC BY 4.0), distributed by arXiv under the Creative Commons Attribution 4.0 International licence, http://creativecommons.org/licenses/by/4.0/, as recorded in the arXiv licence metadata for this exact version and shown on the abstract page https://arxiv.org/abs/2601.06819v1. Full licence notice: \texttt{licenses/arxiv-2601.06819-CC-BY-4.0.txt}.\par\smallskip

\noindent\textit{Editorial scope.} Counted unit: the article's prop prop-up-pry (source0.tex lines 385--402). The article's complete original proof of it is delivered with it (source0.tex lines 404--506, one \texttt{proof} environment of 103 lines). No mathematics is written, repaired or re-derived: the statement and the proof are the article's own lines, in the article's own notation, and no retained line is substituted or re-flowed.  Retained uncounted as the source-local context the proof needs: lines 368--371.  Nothing was omitted from any retained span, so the package carries no omission record.  No retained line contains a citation, so the wrapper carries no bibliography and no bibliography entry.  No retained line contains a figure environment, an image-inclusion command or drawing code, so no image is delivered and the package declares no asset dependency.  Editorial formatting only, in addition: an AMS \texttt{amsart} preamble that restates the article's own declarations and macros used by the retained lines, namely the environments \texttt{lem}, \texttt{prop} on separate counters, together with the standard packages those lines need.  The retained environments print in this document as Lemma 1, Proposition 1 in order of appearance, whereas the article prints the same items with the numbering of its own full document.  The complete native source of the article as distributed in the arXiv e-print for this exact version is bundled unchanged as \texttt{sources/192443/source0.tex}.
\begin{lem}\label{lem-permute-A-B}
Suppose $A'=A \cdot\sigma_1$ and $B'=B\cdot\sigma_2$.
 If $UA =B \cdot\tau$, then $ UA' = B' \cdot (\sigma_2^{-1} \tau \sigma_1) $.
\end{lem}

\begin{prop}\label{prop-up-pry}
Let
\[
C_1 = \begin{pmatrix}
A_1 & A_2 \\
\mathbf{0} & I_{n-d}
\end{pmatrix}
\quad \text{and} \quad
D_1 = \begin{pmatrix}
B_1 & B_2 \\
\mathbf{0} & I_{n-d}
\end{pmatrix},
\]
where $A_1$ and $B_1$ are $d \times d$ matrices, $A_2$ and $B_2$ are arbitrary matrices of appropriate sizes, and $I_{n-d}$ denotes the $(n-d) \times (n-d)$ identity matrix. Then
\begin{equation}\label{equ-up-pyr}
A_1 \simeq_{\mathrm{UP}} B_1 \quad \Longleftrightarrow \quad C_1 \simeq_{\mathrm{UP}} D_1.
\end{equation}
\end{prop}

\begin{proof}
It is clear that
\[
C_1' := \begin{pmatrix}
A_1 & \mathbf{0} \\
\mathbf{0} & I_{n-d}
\end{pmatrix} \simeq_{\mathrm{U}} C_1
\quad \text{and} \quad
D_1' := \begin{pmatrix}
B_1 & \mathbf{0} \\
\mathbf{0} & I_{n-d}
\end{pmatrix} \simeq_{\mathrm{U}} D_1.
\]
Hence, in proving~\eqref{equ-up-pyr}, we may replace $C_1$ and $D_1$ by $C_1'$ and $D_1'$, respectively.

The necessity is straightforward. Suppose there exist $U \in \mathrm{GL}_{d}(\mathbb{Z})$ and $\sigma \in \mathfrak{S}_d$ such that $UA_1 = B_1P_\sigma$. Then the identity
\[
\begin{pmatrix}
U & \mathbf{0} \\
\mathbf{0} & I_{n-d}
\end{pmatrix}
C_1' =
D_1'
\begin{pmatrix}
P_\sigma & \mathbf{0} \\
\mathbf{0} & I_{n-d}
\end{pmatrix}
\]
immediately yields $C_1^{\prime} \simeq_{\mathrm{UP}} D_1^{\prime}$.

We now prove the sufficiency. Suppose there exist $V \in \mathrm{GL}_{n}(\mathbb{Z})$ and $\tau \in \mathfrak{S}_n$ such that $VC_1^{\prime} = D_1^{\prime}P_\tau$. Write the matrices in column form as
\begin{align*}
C_1' &= \big( c_1 \mid \cdots \mid c_{d} \mid \widehat{\mathbf{e}}_1 \mid \cdots \mid \widehat{\mathbf{e}}_{n-d} \big), \\
D_1' &= \big( f_1 \mid \cdots \mid f_{d} \mid \widehat{\mathbf{e}}_1 \mid \cdots \mid \widehat{\mathbf{e}}_{n-d} \big),
\end{align*}
where, for $1 \le i \le d$, the columns $c_i$ and $f_i$ are nonzero only in their first $d$ entries.

By Lemma~\ref{lem-permute-A-B}, we may assume (after a suitable permutation of columns) that
\begin{align}
V c_i &=
\begin{cases}
f_i, & \text{for } 1 \le i \le k, \\
\widetilde{\mathbf{e}}_{i}, & \text{for } k+1 \le i \le d,
\end{cases}
\label{eq:combined1} \\
V \widehat{\mathbf{e}}_i &= f_{i+k}, \quad \text{for all } 1 \le i \le d-k, \label{eq:combined2}
\end{align}
where $k$ is an integer with $0 \le k \le d$.


From the above, we deduce the following unimodular equivalences:
\begin{align}
\big( c_1 \mid \cdots \mid c_k \mid c_{k+1} \mid \cdots \mid c_{d} \big) &\simeq_{\mathrm{U}} \big( f_1 \mid \cdots \mid f_k \mid \widetilde{\mathbf{e}}_{1} \mid \cdots \mid \widetilde{\mathbf{e}}_{d-k} \big), \label{e-c2f} \\
\big( c_1 \mid \cdots \mid c_k \mid \widehat{\mathbf{e}}_{1} \mid \cdots \mid \widehat{\mathbf{e}}_{d-k} \big) &\simeq_{\mathrm{U}} \big( f_1 \mid \cdots \mid f_k \mid \cdots \mid f_{d} \big). \label{e-f2c}
\end{align}

When $k = d$, we may assume without loss of generality that
\[
K \begin{pmatrix} A_1 \\ \mathbf{0} \end{pmatrix} = \begin{pmatrix} B_1 \\ \mathbf{0} \end{pmatrix} P_\sigma,
\]
where
\[
K = \begin{pmatrix} K_{11} & K_{12} \\ K_{21} & K_{22} \end{pmatrix} \in \mathrm{GL}_n(\mathbb{Z}), \quad \sigma \in \mathfrak{S}_d,
\]
with $K_{11} \in \mathbb{Z}^{d \times d}$ and $K_{22} \in \mathbb{Z}^{(n-d) \times (n-d)}$.
It follows that
\[
\begin{pmatrix} K_{11}A_1 \\ K_{21}A_1 \end{pmatrix} = \begin{pmatrix} B_1P_\sigma \\ \mathbf{0} \end{pmatrix}.
\]
Since $A_1$ is nonsingular, we deduce that $K_{21} = \mathbf{0}$. Therefore,
\[
\det(K) = \det(K_{11})\det(K_{22}) = \pm 1,
\]
which implies $K_{11} \in \mathrm{GL}_d(\mathbb{Z})$. From the identity $K_{11}A_1 = B_1P_\sigma$, we conclude that $A_1 \simeq_{\mathrm{UP}} B_1$.


Now suppose $k < d$. From \eqref{e-c2f}, we first derive the equivalence
\begin{align}
\label{e-c2cp}
(c_1 \mid \cdots \mid c_k \mid c_{k+1} \mid \cdots \mid c_d) \simeq_{\mathrm{U}} (c_1 \mid \cdots \mid c_k \mid \widehat{\mathbf{e}}_1 \mid \cdots \mid \widehat{\mathbf{e}}_{d-k}).
\end{align}

Note that the last $n - d$ components of each column vector $c_j$ and $f_j$ (for $j = 1, \dots, d$) are zero. Hence, for each $i = k+1, \dots, d$, the row vector $\mathbf{e}_i^\mathsf{T}$ can be expressed as an integer linear combination of the first \(d\) rows of the matrix \((c_1 \mid \cdots \mid c_k \mid \cdots \mid c_{d})\).

By inserting $\mathbf{e}_i^\mathsf{T}$ as a row into the matrix $(c_1 \mid \cdots \mid c_d)$, we obtain
\[
(c_1 \mid \cdots \mid c_k \mid c_{k+1} + \widehat{\mathbf{e}}_1 \mid \cdots \mid c_d + \widehat{\mathbf{e}}_{d-k}).
\]
Since $\widehat{\mathbf{e}}_i \neq \widehat{\mathbf{e}}_j$ for $i \neq j$, and since each $c_{k+i} + \widehat{\mathbf{e}}_i$ can be transformed into $\widehat{\mathbf{e}}_{k+1}$ via elementary row operations, we obtain
\begin{align*}
(c_1 \mid \cdots \mid c_d) &\simeq_{\mathrm{U}} (c_1 \mid \cdots \mid c_k \mid c_{k+1} + \widehat{\mathbf{e}}_1 \mid \cdots \mid c_d + \widehat{\mathbf{e}}_{d-k}) \\
&\simeq_{\mathrm{U}} (c_1 \mid \cdots \mid c_k \mid \widehat{\mathbf{e}}_1 \mid \cdots \mid \widehat{\mathbf{e}}_{d-k}),
\end{align*}
as required.

Combining the equivalences \eqref{e-f2c} and \eqref{e-c2cp}, we conclude that
\begin{align}
\label{e-cf}
\begin{pmatrix} A_1 \\ \mathbf{0} \end{pmatrix}
= (c_1 \mid \cdots \mid c_d) &\simeq_{\mathrm{UP}} (f_1 \mid \cdots \mid f_d) = \begin{pmatrix} B_1 \\ \mathbf{0} \end{pmatrix}.
\end{align}
By an argument analogous to the case $k = d$, it follows that $A_1 \simeq_{\mathrm{UP}} B_1$.
\end{proof}

\end{document}
